\documentclass{beamer}
\usepackage{pifont}
\input{unicode-shim.tex}
\newbox\unicodeprobe
\newbox\controlprobe

\begin{document}
% Compile twice and inspect the actual PDF outline title: it must remain α → β.
\section{α → β}
\begin{frame}{Unicode workflow probes}
% A pasted bullet must retain the same math glyph and spacing as \bullet.
\setbox\unicodeprobe=\hbox{$a • b$}
\setbox\controlprobe=\hbox{$a \bullet b$}
\typeout{PPS-BULLET-UNICODE-WIDTH=\the\wd\unicodeprobe}
\typeout{PPS-BULLET-CONTROL-WIDTH=\the\wd\controlprobe}
\ifdim\wd\unicodeprobe=\wd\controlprobe\else
  \errmessage{Unicode bullet width differs from its math control}
\fi
Unicode bullet: $a • b$; control: $a \bullet b$.

% Permit breaks after ordinary hyphens, but disable word hyphenation so this
% probe isolates the pasted U+2011 non-breaking hyphen.
\begingroup
\hyphenpenalty=10000 \exhyphenpenalty=0
\setbox\unicodeprobe=\vbox{\hsize=35pt\rightskip=0pt plus1fil
  \noindent state‑of‑the‑art\par}
\setbox\controlprobe=\vbox{\hsize=35pt\rightskip=0pt plus1fil
  \noindent state\mbox{-}of\mbox{-}the\mbox{-}art\par}
\typeout{PPS-NBH-UNICODE-HEIGHT=\the\ht\unicodeprobe}
\typeout{PPS-NBH-CONTROL-HEIGHT=\the\ht\controlprobe}
\ifdim\ht\unicodeprobe=\ht\controlprobe\else
  \errmessage{Unicode non-breaking hyphen wraps differently from its control}
\fi
\par Unicode non-breaking hyphen:\par\box\unicodeprobe
\par No-break control:\par\box\controlprobe
\endgroup
\end{frame}
\end{document}
